\subsection{What a two-generator ordinary comparison would require}

The order-five result bounds the original rational strict group by
one generator.  On the split branches of
KT-RESIDUAL-CYCLICITY-SCOPE, a nonzero rational group and the
separately constructed negative-conjugation class give
\[
 \dim_k\mathcal H_E\geq2.                                      \tag{D2.1}
\]
There is no upper bound of two for this entire ordinary group in
the results proved here.  In particular, ``dimension two'' below
is a proposed additional hypothesis, not a consequence of the
quintic class-number estimate.

\paragraph{Published comparisons do allow more than one extension.}
Berger--Klosin, \emph{On lifting and modularity of reducible residual
Galois representations over imaginary quadratic fields}, IMRN 2015,
use no dimension-one assumption for their extension space.
In the accepted author version, arXiv:1606.06535, Theorem~8.2
(pp.~27--28) requires a projectively unique modular basis and
finitely many characteristic-zero lifts for each member, together
with the preceding local hypotheses and a divisible Selmer upper
bound matched by a congruence lower bound.  Theorem~8.5 (p.~28)
replaces projective uniqueness by
\[
 \varpi H^1_\Sigma
       (E,\Psi^{-1}\otimes\operatorname{Frac}(\mathcal O)/\mathcal O)
       =0.                                                     \tag{D2.2}
\]
It retains the other hypotheses.  Assumptions~2.5 and~4.2
(pp.~6, 12) explicitly require
\[
 \#H^1_\Sigma(E,\Psi^{-1}\otimes
              \operatorname{Frac}(\mathcal O)/\mathcal O)
       \leq\#\mathcal O/(L)\leq\#T/J.                          \tag{D2.3}
\]
Lemma~4.4 (p.~13) uses Fitting ideals to compare the lattice
quotient with this Selmer group.  The final surjectivity in that
lemma uses both inequalities in (D2.3); it does not furnish the
left inequality from modularity.

These are two-constituent rank-two results.  The imaginary
quadratic principality argument, Proposition~7.1 and its proof
(p.~23), assumes an anticyclotomic residual ratio
$\bar\Psi^c=\bar\Psi^{-1}$.  For the Tate ratio of order five,
\[
 (\omega^{e}|_{G_E})^c=\omega^{e}|_{G_E},
 \qquad \omega^{e}|_{G_E}\ne\omega^{-e}|_{G_E},                 \tag{D2.4}
\]
provided $E$ is prime to $p$, as in the auxiliary-field setup.
Indeed $E\cap\mathbf Q(\mu_p)=\mathbf Q$ here, so restriction
preserves the order five.  A common twist does not change a
constituent ratio, and interchanging the two constituents only
inverts it.  Thus neither operation repairs (D2.4).
The rank-three residual representation also cannot be treated as
two absolutely irreducible constituents by combining two of its
three distinct characters into a block.

The published conclusion concerns the characteristic-zero
quotient $R_\tau^0$: the authors explicitly separate it from
possible characteristic-$p$ components, in Section~2.5
(pp.~8--9).  It supplies no injection of our integral odd
off-diagonal module into the resulting trace algebra.

\paragraph{The divisible-by-$p$ weight does not verify (D2.2).}
KT-ORDER-FIVE-PDIV-WEIGHT proves
\[
 p\,\mathcal S_p(e)=0\qquad(p\mid e),                         \tag{D2.5}
\]
for the rational strict group with coefficient
$(\mathbf Q_p/\mathbf Z_p)(\chi_p^{-e})$.
Restriction places these classes in the positive-conjugation
part of the corresponding group over $E$; it does not identify
them with its whole divisible Selmer group.  It bounds neither
the negative-conjugation part nor a relaxation at any allowed
ramified prime.  Moreover the character in (D2.2) is the
anticyclotomic character of the cited theorem, and its local
condition is the theorem's crystalline condition at every
prime over $p$ (Section~2.3, pp.~5--6).  These are not supplied
by the inclusion of the original rational strict group.
Even with the coefficient ring $\mathbf Z_p$, (D2.5) therefore
does not verify (D2.2).  With a ramified coefficient extension,
annihilation by $p$ would additionally be weaker than
annihilation by its uniformizer.

There is a separate local obstruction to applying that paper to
the actual large-weight Tate lifts: its ``crystalline'' deformation
condition is the short Fontaine--Laffaille condition
$\operatorname{Fil}^0=D$, $\operatorname{Fil}^{p-1}=0$
(Section~2.1, pp.~3--4).  If $p\mid e$ and $e>0$ is even, then
$e\geq2p$.  The Hodge--Tate weight difference $e$ of an adjacent
Tate pair cannot be put in that interval by a common twist.
Replacing $e$ by a congruent smaller integer changes the
characteristic-zero local condition and the divisible
coefficient group, so it does not prove the desired application.

\paragraph{Nonprincipal Eisenstein ideals are compatible with $R=T$.}
Wake--Wang-Erickson, \emph{The Eisenstein ideal with squarefree
level}, author version dated 19 January 2021,
Theorem~6.4.1 (p.~33), proves an actual $R=T$ theorem with
\[
 I/I^2\simeq
 \mathbf Z_p/(\ell_0-1)\oplus\mathbf Z_p/(\ell_1-1).            \tag{D2.6}
\]
Its hypotheses are $N=\ell_0\ell_1$, both
$\ell_i\equiv1\pmod p$, Atkin--Lehner signs $(-1,-1)$, and
neither prime being a $p$th power modulo the other.  Its
deformation problem is the rank-two residual
pseudorepresentation $1\oplus\omega$, with determinant
$\chi_p$ and the weight-two finite-flat condition at $p$.
Here $p\nmid N$ follows from $\ell_i\equiv1\pmod p$.
The proof computes the conormal module and verifies the
Wiles numerical criterion; Proposition~5.2.1 (p.~30) explains
that numerical comparison.  This supplies a genuine example
where one generator is unnecessary.  It supplies no such
conormal calculation for $1\oplus\omega^e\oplus\omega^{-e}$.
All pairwise ratios in our triple have order five, whereas the
ratio in this theorem is $\omega$, of even order $p-1$; twisting
or permuting the constituents does not identify the problems.

Berger--Klosin, \emph{Modularity of residual Galois extensions
and the Eisenstein ideal}, Trans. AMS 2019, accepted author
version, Proposition~3.9 (p.~7) and Theorem~4.8 (p.~10), also
allow a noncyclic extension space.  Under the stated matched
divisible-Selmer/congruence bound and principality of the
individual residual-branch ideals, nonprincipality of the
global Eisenstein ideal forces \emph{more} modular extension
lines than the dimension of the Selmer group.
Proposition~4.6 (p.~10) proves the needed branch principality
when the opposite extension space is one-dimensional.
This is not a vanishing criterion for either extension space.

\begin{lemma}[KT-GMA-MULTIGENERATOR-CONORMAL]
Let $A$ be a complete Noetherian local ring and let the
three-block generalized matrix algebra, trace $T$, and
anti-involution satisfy Assumption~2.7 of Berger--Klosin,
arXiv:2611.19007v1.  In particular the anti-involution
interchanges blocks one and three, fixes block two, and all
off-diagonal modules $A_{ij}$ are finite over $A$.
Write $A'_{ij}=A_{ik}A_{kj}$ for the remaining index $k$, and
let $I\subset\mathfrak m_A$ be the total reducibility ideal.
Assume
\[
 A_{12}=A'_{12},\qquad
 A_{13}/A'_{13}\text{ is cyclic},\qquad
 A_{32}/A'_{32}\text{ is generated by at most }d\text{ elements}.
                                                               \tag{D2.7}
\]
Put $B=A_{21}/A'_{21}$ and $C=A_{31}/A'_{31}$.
For any $A$-module $M$ for which the right side is finite,
\[
 \#\operatorname{Hom}_A(I/I^2,M)
 \leq\#\operatorname{Hom}_A(C,M)\,
       \#\operatorname{Hom}_A(B,M)^d.                         \tag{D2.8}
\]
If $A/I\simeq\mathcal O$ is a complete discrete valuation
ring with finite residue field and
$\operatorname{Hom}_A(B,\mathcal O)=
 \operatorname{Hom}_A(C,\mathcal O)=0$, then $I/I^2$ is finite.
For every $n$ annihilating it by $\varpi^n$, (D2.8) applies
with $M=\mathcal O/\varpi^n$ and with its left side replaced
by $\#I/I^2$.
\end{lemma}

\begin{proof}
This is a generator-count variant of the complete proof of
Proposition~2.12 and Corollary~2.16 in the cited preprint
(pp.~7--9).  Put
\[
 N=A_{31}/A_{32}A_{23}A_{31},\qquad
 K=A'_{31}/A_{32}A_{23}A_{31}.                                \tag{D2.9}
\]
The product identities of Lemma~2.9, Remark~2.10, and
Lemma~2.11, using
$A_{12}=A'_{12}$, give $I=T(A_{13}A_{31})$.
Multiplication gives the surjections
\[
 (A_{13}/A'_{13})\otimes_A N\longtwoheadrightarrow I/I^2,
 \qquad
 (A_{32}/A'_{32})\otimes_A B\longtwoheadrightarrow K.          \tag{D2.10}
\]
These are precisely the well-defined product maps (2.7) and
(2.10) of that proof; their construction does not require
cyclicity of $A_{32}/A'_{32}$.
Choose a cyclic generator for the first factor of the first
map, and at most $d$ generators for the first factor of the
second.  We obtain $N\twoheadrightarrow I/I^2$ and
$B^{\oplus d}\twoheadrightarrow K$.
The definition gives $0\to K\to N\to C\to0$.
Applying $\operatorname{Hom}_A(-,M)$ therefore gives an
injection from $\operatorname{Hom}_A(I/I^2,M)$ into
$\operatorname{Hom}_A(N,M)$ and an exact sequence
\[
 0\longrightarrow\operatorname{Hom}_A(C,M)
 \longrightarrow\operatorname{Hom}_A(N,M)
 \longrightarrow\operatorname{Hom}_A(B,M)^{\oplus d}.          \tag{D2.11}
\]
Exactness is asserted at the first two terms; the last arrow
need not be surjective.  This proves (D2.8).
Taking $M=\mathcal O$ proves that
$\operatorname{Hom}_{\mathcal O}(I/I^2,\mathcal O)=0$.
Since $I/I^2$ is finite over $A/I=\mathcal O$, it is torsion,
and hence finite.  The elementary divisor theorem identifies
its cardinality with that of its dual into
$\mathcal O/\varpi^n$ whenever $\varpi^n$ annihilates it.
\end{proof}

\paragraph{The exact extra arithmetic inequality.}
Suppose, conditionally, that the full deformation and Hecke
setup of the cited three-character preprint is satisfied,
except that the residual Selmer bound in (S2)(c) is replaced
by dimension at most $d$.  Retain (S2)(a), (S2)(b), the
finite divisible groups in (S1), the ordinary GMA injections
of Theorem~6.4, and all the Hecke and local hypotheses.
The preceding lemma replaces the conormal estimate in the
proof of Theorem~10.2 by
\[
 \operatorname{length}_{\mathcal O}(I/I^2)
 \leq\operatorname{length}_{\mathcal O}(H_{\rm long})
       +d\operatorname{length}_{\mathcal O}(H_{\rm adj}),       \tag{D2.12}
\]
where, in the notation of that theorem, and with
$V=\operatorname{Frac}(\mathcal O)/\mathcal O$,
\[
 \begin{aligned}
 H_{\rm adj}&=H^1_{\rm Gr}
   (G_\Sigma,V(\Psi_k^{-1}\chi_p)),\\
 H_{\rm long}&=H^1_{\rm Gr}
   (G_\Sigma,V(\Psi_k^{-1-c}\chi_p^2))^{c=-1}.
 \end{aligned}                                                \tag{D2.13}
\]
Consequently its Theorem~7.8 still yields $R=T^0\times_{T^0/J^0}
\mathcal O$ if one independently proves
\[
 \operatorname{length}_{\mathcal O}(H_{\rm long})
       +d\operatorname{length}_{\mathcal O}(H_{\rm adj})
 \leq\operatorname{length}_{\mathcal O}(T^0/J^0)<\infty.       \tag{D2.14}
\]
This is a conditional numerical criterion, not an application
to the actual KTOWER source.  For $d=2$ the adjacent Selmer
length is counted twice.  The original proof's two $L$-value
factors bound one copy of each divisible group; they do not
automatically establish (D2.14) with $d=2$.
Even a residual upper bound of two would not bound either
integral length or compute the congruence module.

\paragraph{Remaining boundary.}
For the present Tate triple we have not established the
source's ordinary-flag identification, its other residual
Selmer assumptions, a two-generator bound for the entire
adjacent ordinary group, or (D2.14).  Nor has a saturated
integral comparison been proved between the original odd
off-diagonal module and a trace algebra.  The order-five
class-component estimate and its $p\mid e$ consequence do
not supply these missing statements.  Thus this bounded
audit gives an explicit replacement inequality and records
genuine noncyclic comparison theorems, but removes no
remaining prime, character, or source point.

\paragraph{Source and review scope.}
The new algebraic proof above was checked against the whole
proof of Berger--Klosin 2026 Proposition~2.12, Corollary~2.16,
Theorem~7.8, and Theorem~10.2.  The cited 2015 proofs of
Lemma~4.4, Theorem~6.3, Proposition~7.1, and
Theorems~8.2/8.5, and the stated setup and assumptions, were
read; their imported modularity, lattice, and numerical
theorems remain published inputs.  The complete short
proofs of Wake--Wang-Erickson Proposition~5.2.1 and
Theorem~6.4.1 and Berger--Klosin 2019 Proposition~3.9 and
Theorem~4.8 were read for scope.  No claim is made to
reprove those papers or to have kernel-checked their
arguments.  The primary PDFs, exact retrieval hashes, and
audit records are retained separately.
